\section{個案分析：The DoDo Men－嘟嘟人}
\label{sec:case}

本章依第~\ref{sec:method} 節的分析框架，逐一回答四個研究問題。除特別標明外，
結構與互動密度指標以主留言為口徑，情緒指標以含回覆的語意標註為口徑（見表~\ref{tab:scope}）。

\subsection{頻道整體社群健康（RQ1）}
\label{sec:case-health}

\textbf{規模與分層。} The DoDo Men 在分析範圍內共 351 支影片、247{,}068 則主留言、112{,}108 位留言者，
平均每支影片約 704 則主留言。表~\ref{tab:tiers} 呈現以涵蓋率分層的結果：一次性留言者雖占 68.1\%，但其
benchmark 百分位僅 PR~15（明顯\textbf{低於} cohort 平均 78.5\%），代表 The DoDo Men 的留言比同儕\textbf{更
不依賴}一次性流量。占比僅 3.8\% 的核心與常客層（涵蓋率 $\ge 2\%$）貢獻了 29.3\% 的留言（PR~56），其中
僅 990 位核心留言者（0.9\%）平均留言過 32.5 支影片，貢獻 14.2\% 的留言，顯示存在一群高黏著的主力觀眾。

\begin{table}[H]
\centering
\caption{留言者涵蓋率分層（主留言口徑）}
\label{tab:tiers}
\small
\begin{tabular}{lrrrr}
\toprule
\textbf{層級} & \textbf{占留言者} & \textbf{留言數} & \textbf{留言貢獻} & \textbf{人均影片}\\
\midrule
core（核心，涵蓋率 $\ge 5\%$） & 0.9\% & 35{,}063 & 14.2\% & 32.5 \\
regular（常客，$2$–$5\%$） & 2.9\% & 37{,}207 & 15.1\% & 10.9 \\
returning（回訪，$\ge 2$ 支） & 28.1\% & 97{,}289 & 39.4\% & 3.0 \\
one-time（一次性，1 支） & 68.1\% & 77{,}509 & 31.4\% & 1.0 \\
\bottomrule
\end{tabular}
\end{table}

\textbf{回訪、情緒與相對基準。} 表~\ref{tab:health-bench} 將關鍵指標對照 48 頻道 cohort。The DoDo Men 的
跨影片回訪率 24.9\%（PR~81）與近期 rolling 回訪率 30.5\%（PR~91）均明顯高於同儕，顯示留言者會隨後續影片
持續參與。情緒面，正面率 47.4\%（PR~73）偏高、負面率 7.5\%（PR~38）與按讚加權負面率 5.2\%（PR~35）皆\textbf{低
於} cohort 平均，整體情緒健康。

\begin{table}[H]
\centering
\caption{頻道健康關鍵指標對照 48 頻道 cohort（PR 為百分位）}
\label{tab:health-bench}
\small
\begin{tabular}{lrrr}
\toprule
\textbf{指標} & \textbf{本頻道} & \textbf{cohort 平均} & \textbf{PR}\\
\midrule
跨影片回訪率（4 影片窗） & 24.9\% & 16.7\% & 81 \\
近期回訪率（rolling 50/50） & 30.5\% & 17.6\% & 91 \\
正面留言率 & 47.4\% & 35.7\% & 73 \\
負面留言率 & 7.5\% & 11.8\% & 38 \\
按讚加權負面率 & 5.2\% & 11.3\% & 35 \\
每千次觀看主留言數 & 0.80 & 1.13 & 40 \\
一次性留言者占比 & 68.1\% & 78.5\% & 15 \\
\bottomrule
\end{tabular}
\end{table}

\textbf{社群結構。} co-commenter network（19{,}216 nodes、1{,}796{,}606 edges）經 Louvain 偵測出
\textbf{3 個觀眾社群}，modularity~$=0.31$、HHI~$=0.37$、最大社群占 47.7\%、effective communities~$=2.69$。
其中 modularity~$0.31$ 雖屬絕對值的中等水準（觀眾天生重疊），但相對 cohort（中位數約 0.22）偏高，代表
The DoDo Men 的留言者比多數頻道更能被分出偏好傾向；不過社群高度集中（社群集中度指數 87／100，非常高）。

\textbf{綜合判讀。} 整合上述指標，本研究將 The DoDo Men 歸類為「\textbf{穩定討論引擎型}」頻道：互動轉換力
（70）與核心黏著度（67）偏高、情緒風險壓力（48）中段、但社群集中度（87）非常高（表~\ref{tab:indices}）。
換言之，頻道擁有穩定且願意回訪的主力留言社群、整體情緒健康，但互動高度集中於少數觀眾群，內容擴張時需留意是否
能觸及核心以外的受眾。\textbf{對 RQ1 的回答}：The DoDo Men 的留言社群整體健康，互動主要來自穩定核心與回訪
觀眾，而非一次性流量。

\begin{table}[H]
\centering
\caption{頻道社群健康六大指數（相對 48 頻道 cohort，0--100）}
\label{tab:indices}
\small
\begin{tabular}{lcc}
\toprule
\textbf{指數} & \textbf{分數} & \textbf{相對水準}\\
\midrule
互動轉換力 & 70 & 偏高 \\
核心黏著度 & 67 & 偏高 \\
討論品質 & 64 & 中段 \\
情緒風險壓力 & 48 & 中段 \\
社群集中度 & 87 & 非常高 \\
內容組合連通性 & 51 & 中段 \\
\bottomrule
\end{tabular}
\end{table}

\subsection{觀眾社群分析（RQ2）}
\label{sec:case-community}

The DoDo Men 的留言社群並非「一群粉絲」，而是由三個具不同內容偏好與情緒特徵的觀眾群組成
（表~\ref{tab:personas}）。需提醒，由於 modularity 僅 0.31，這些社群應理解為\textbf{軟性偏好傾向}而非
壁壘分明的派系，成員之間存在大量重疊。

\begin{table}[H]
\centering
\caption{三個觀眾社群（audience segment profile）}
\label{tab:personas}
\small
\begin{tabularx}{\linewidth}{l l Y c}
\toprule
\textbf{命名} & \textbf{規模} & \textbf{相對偏好主題（affinity lift）} & \textbf{正／負面率}\\
\midrule
知識職涯型 & 47.7\%（9{,}175） & 教育建議 $1.68\times$、職場科技 $1.45\times$ & 50.7\%／5.1\% \\
挑戰冒險型 & 32.0\%（6{,}148） & 荒野求生 $2.43\times$、爭議回應 $2.14\times$、品牌合作 $1.67\times$ & 47.2\%／8.1\% \\
旅遊來賓型 & 20.3\%（3{,}893） & 旅遊／來賓為主，無明顯 over-index & 50.7\%／6.6\% \\
\bottomrule
\end{tabularx}
\end{table}

最大的\textbf{知識職涯型}（47.7\%）相對偏好教育與職場科技內容，且負面率最低（5.1\%），是頻道最穩定的觀眾群；
\textbf{挑戰冒險型}（32.0\%）對荒野求生、爭議回應與品牌合作內容特別投入，但負面率最高（8.1\%），屬高互動但
較敏感的觀眾群；\textbf{旅遊來賓型}（20.3\%）則為旅遊與來賓內容的基底觀眾。\textbf{對 RQ2 的回答}：頻道確實
存在三個內容偏好與情緒特徵不同的觀眾社群。

\subsection{觀眾社群與內容敏感度（RQ3 之一）}
\label{sec:case-sensitivity}

進一步在\textbf{同一主題內}比較不同社群的負面率（控制題材），可辨識「多群都看、但某群特別容易負面」的內容敏感度
差異（表~\ref{tab:sensitivity}）。例如三群都會留言「職場科技」內容，但挑戰冒險型的負面率（14\%）明顯高於
知識職涯型（5\%），差距達 8 個百分點；「產品開箱」則以知識職涯型最易負面（12\%）。相對地，「荒野求生」三群
負面率接近（$\sim$10--11\%），代表此主題的反應在各群間一致，差異主要由主題本身而非特定社群造成。

\begin{table}[H]
\centering
\caption{同一主題內不同社群的負面率（僅列各群差距 $\ge 3$pp 之主題）}
\label{tab:sensitivity}
\small
\begin{tabular}{lcccc}
\toprule
\textbf{主題} & \textbf{知識職涯型} & \textbf{挑戰冒險型} & \textbf{旅遊來賓型} & \textbf{群間差距}\\
\midrule
職場科技 & 5\% & \textbf{14\%} & 9\% & 8pp \\
產品開箱 & \textbf{12\%} & 7\% & 10\% & 5pp \\
團隊日常 & 3\% & \textbf{7\%} & 6\% & 4pp \\
\midrule
荒野求生（對照：無顯著差異） & 11\% & 10\% & 10\% & $\sim$0 \\
\bottomrule
\end{tabular}
\end{table}

此分析直接支援分眾內容策略與風險監控：面向挑戰冒險型的職場科技／團隊日常內容、以及面向知識職涯型的產品開箱
內容，發布後應提高留言監控強度。

\subsection{情緒風險、語意主體與留言衝突（RQ3 之二、RQ4）}
\label{sec:case-risk}

\textbf{被讚與被罵的面向（ABSA）。} 全頻道層級，觀眾正面情緒主要針對\textbf{主持人}（host persona，35\%）
與\textbf{內容品質}（28\%），顯示人格魅力與內容是頻道的正面資產；負面情緒則主要針對\textbf{價值觀／政治}
（28\%）與\textbf{真實性／造假}（28\%），其次為內容品質（19\%）與主持人（16\%）。這說明頻道的負面風險較高層級
（聲譽／價值觀）而非單純內容優化問題。

\textbf{負面熱點影片。} 表~\ref{tab:neg-videos} 列出負面率與按讚加權負面率較高的影片。值得注意的是「一天內用
十元硬幣換到兩張機票」raw 負面率 28.4\%、但\textbf{按讚加權負面率高達 49.3\%}，代表該片的負面意見被其他觀眾
大量按讚放大，擴散風險高於其表面負面率。

\begin{table}[H]
\centering
\caption{負面熱點影片（依按讚加權負面率）}
\label{tab:neg-videos}
\small
\begin{tabularx}{\linewidth}{Y rrr}
\toprule
\textbf{影片} & \textbf{負面率} & \textbf{按讚加權負面} & \textbf{留言數}\\
\midrule
一天內用十元硬幣換到兩張機票 & 28.4\% & 49.3\% & 1{,}105 \\
帶 Enzo 體驗最道地的行程！算命超準？！ & 34.0\% & 44.6\% & 1{,}552 \\
跟倉儲職業級高手的開箱合作 & 35.8\% & 30.8\% & 1{,}486 \\
蘋果 vs 三星｜開箱包裝體驗差超多 & 23.4\% & 27.3\% & 2{,}059 \\
\bottomrule
\end{tabularx}
\end{table}

\textbf{留言衝突熱點（RQ4）。} 將「負面多」與「真的吵起來」區分後，表~\ref{tab:conflict-videos} 列出衝突
熱點影片。多數高衝突影片以\textbf{對立母串}為主、\textbf{圍剿}少（全頻道對立母串 833 串、圍剿僅 114 串，兩者
相關約 0.6），代表衝突多為「回覆與母留言立場相反的爭論」而非「一面倒的負面圍攻」。少數例外如「到台灣高中第一
志願教課」同時具高圍剿（8）與高對立母串（16），屬需優先關注的留言串。\textbf{對 RQ4 的回答}：負面與衝突並非
同一件事；「巴拉圭邦交國」「荒島求生 Day1」「倉儲開箱」「建中教課」等影片不只是負評多，而是引發了實質的留言
互動衝突。

\begin{table}[H]
\centering
\caption{留言衝突熱點影片（依 like-weighted conflict score）}
\label{tab:conflict-videos}
\small
\begin{tabularx}{\linewidth}{Y cccc}
\toprule
\textbf{影片} & \textbf{衝突串} & \textbf{圍剿} & \textbf{對立母串} & \textbf{score}\\
\midrule
巴拉圭人超愛台灣！開箱邦交國外交生活 & 9 & 0 & 8 & 6.59 \\
荒島求生 Day1：自己 DIY 蓋避難所 & 16 & 0 & 17 & 6.53 \\
跟倉儲職業級高手的開箱合作 & 16 & 1 & 14 & 6.43 \\
到台灣高中第一志願教課（建中） & 17 & \textbf{8} & 16 & 5.17 \\
從荷蘭帶小帥哥回台灣 & 20 & 3 & 15 & 4.84 \\
\bottomrule
\end{tabularx}
\end{table}

\textbf{對 RQ3 的回答}：正面主要針對主持人與內容品質，負面主要針對價值觀／政治與真實性；少數影片的負面被按讚
放大，屬較高層級的聲譽風險而非單純內容問題。

\subsection{策略建議}
\label{sec:case-strategy}

\textbf{內容策略。} 核心粉絲（知識職涯型、旅遊來賓型）偏好的教育、職場、旅遊與來賓內容情緒穩定、負面低，適合
作為維持社群黏著的主力，並可系統化為可複製的企劃公式；挑戰冒險型偏好的荒野求生、爭議回應與品牌合作雖帶來高
互動，但負面與衝突風險較高，需搭配更清楚的表達與風險監控。

\textbf{分眾溝通策略。} 不同社群需要不同溝通方式：知識職涯型重視清楚資訊與證據；挑戰冒險型對價值觀、真實性與
事件回應較敏感，需要更透明、即時且一致的回應。面向挑戰冒險型的職場科技內容、面向知識職涯型的產品開箱內容，是
群間負面差異最大的組合，應優先準備溝通與說明素材。

\textbf{炎上與風險監測策略。} 頻道不應只看負面率，而應同時監控：(1) 按讚加權負面率（負面是否被放大）、
(2) conflict score 與圍剿／對立母串數、(3) 新留言者負面率、(4) 各觀眾社群的情緒變化。
若多個風險指標同時升高（如負面被放大、衝突上升、特定社群情緒轉差），應優先處理。本研究的個案顯示，The DoDo Men
整體情緒健康，真正需要優先關注的是少數「負面被按讚放大」與「引發對立母串」的影片，而非整體負面水準。
